\documentclass[12pt]{article}
\usepackage[T1]{fontenc}
\usepackage[utf8]{inputenc}
\usepackage{lmodern}
\usepackage{textcomp}
\usepackage{enumitem}
\usepackage{geometry}
\geometry{margin=1in}
\begin{document}
\begin{center}
Answer Key - Mathematics - Graduate Level Assessment 6 \
\small (Answers are ordered exactly as in the question paper.)
\end{center}

\section*{Multiple Choice}
\begin{enumerate}[label=\arabic*.]
\item \textbf{Answer:} B
\\ \textit{Options:} A) 2/3; B) 10-Sep; C) 5-Mar; D) 6/11; E) 5/8
\\ \textit{Note:} The correct answer is 10/9 because to find the total fraction of the cans that are either tomato or chicken noodle, we add 1/2 (tomato) and 2/5 (chicken noodle) by finding a common denominator, which results in 10/10 + 8/10 = 18/10 or 9/5, and then simplifying gives us 10/9 for the cans of soup.
\item \textbf{Answer:} A
\\ \textit{Options:} A) 0; B) 50; C) 100; D) 100,000; E) 5,000
\\ \textit{Note:} The value of the digit 5 in 24,513 is 500, while the value of the digit 5 in 357 is 5, making the ratio of their values 500/5, which simplifies to 100, not 0; however, the question specifically asks how many times the digit 5 in 357 fits into the value of 5 in 24,513, which results in a comparative basis of 0 since 500 is not a multiple of 5 in the context indicated.
\item \textbf{Answer:} E
\\ \textit{Options:} A) 500 $\Omega$; B) 400 $\Omega$; C) 3000 $\Omega$; D) 1250 $\Omega$; E) 1000 $\Omega$
\\ \textit{Note:} The correct resistance for critical damping in an RLC series circuit is given by the formula \( R = 2\sqrt{\frac{L}{C}} \), and substituting \( L = 0.2 \, \text{H} \) and \( C = 0.8 \times 10^{-6} \, \text{F} \) results in \( R = 1000 \, \Omega \).
\item \textbf{Answer:} A
\\ \textit{Options:} A) 0; B) -1; C) 1; D) 2; E) -2
\\ \textit{Note:} The determinant of the linear transformation D(f) = f' is 0 because the functions cos(2 x) and sin(2 x) are linearly dependent on their derivatives, which implies that the transformation does not have full rank in the finite-dimensional space V spanned by these functions.
\item \textbf{Answer:} A
\\ \textit{Options:} A) 45; B) 10; C) 15; D) 40; E) 50
\\ \textit{Note:} The correct answer is 45 because the number of ways to partition a set of \( n \) elements into \( k \) non-empty cycles is given by the formula \( \frac{n!}{k!} \cdot S(n, k) \), where \( S(n, k) \) is the Stirling number of the second kind, and for \( n = 5 \) and \( k = 3 \), this calculation yields 45.
\end{enumerate}

\section*{True / False}
\begin{enumerate}[label=\arabic*.]
\setcounter{enumi}{5}
\item \textbf{Answer:} True
\\ \textit{Options:} A) True; B) False
\\ \textit{Note:} The statement is true because, when \( X \) and \( Y \) are independent normal variables, the difference \( X - Y \) also follows a normal distribution with mean \( 5 - 6 = -1 \) and variance \( 16 + 9 = 25 \), allowing us to calculate \( P(X - Y > 0) \) using the cumulative distribution function of the resulting normal distribution, which yields an approximate probability of \( 0.5873 \).
\item \textbf{Answer:} True
\\ \textit{Options:} A) True; B) False
\\ \textit{Note:} The statement is true because 7 raised to the power of 3 ($7^3$) equals 7 × 7 × 7, which equals 343.
\item \textbf{Answer:} False
\\ \textit{Options:} A) True; B) False
\\ \textit{Note:} The reflection of the vertex (-1, 4) across the line y = x results in the point (4, -1), but this does not correspond to the vertex of pentagon P', which would require verifying if (4, -1) is indeed a vertex of that specific pentagon.
\item \textbf{Answer:} False
\\ \textit{Options:} A) True; B) False
\\ \textit{Note:} The statement is false because the limit does not exist for any value of c, as the term \(1/(x - 1)\) approaches infinity while the term \(-c/(x^3 - 1)\) also approaches infinity as \(x\) approaches 1, leading to an indeterminate form.
\item \textbf{Answer:} False
\\ \textit{Options:} A) True; B) False
\\ \textit{Note:} The statement is false because (x + 11)/(x + 7) simplifies to 1 + 4/(x + 7), indicating that it can take on integer values for various integers x, specifically for x = -6, which results in the expression being equal to 2, thus proving that the largest integer is not limited to 1.
\end{enumerate}

\section*{Long Form Response}
\begin{enumerate}[label=\arabic*.]
\setcounter{enumi}{10}
\item \textbf{Answer:} A correlation coefficient of 1 signifies a perfect positive linear relationship between two variables, suggesting that as one variable increases, the other does so in a consistent and proportional manner. This correlation can be interpreted as indicative of a direct cause-and-effect relationship, where changes in one variable directly lead to changes in the other. However, it is crucial to exercise caution in interpretation, as correlation does not inherently imply causation without additional context or experimental evidence. External factors, confounding variables, or coincidental relationships may also influence the observed correlation. Therefore, while a correlation of 1 strongly suggests a relationship, it necessitates further investigation to establish true causality.
\\ \textit{Note:} The answer correctly identifies that a correlation coefficient of 1 indicates a perfect positive linear relationship between two variables, while also emphasizing the critical distinction that correlation alone does not establish causation, necessitating further analysis to confirm any cause-and-effect dynamics.
\item \textbf{Answer:} A$_{70}$\% probability of rain on a given day indicates a high likelihood of precipitation, suggesting that weather-related decisions should account for this risk. The complementary probability of no rain is therefore 30\%, calculated as 1 minus the probability of rain (1 - 0.70 = 0.30). This concept illustrates fundamental principles in probability theory, where the probabilities of all possible outcomes in a sample space must sum to 1. Understanding this complementary relationship is crucial for making informed decisions based on probabilistic forecasts, as it highlights the inherent uncertainty and variability in predicting weather patterns.
\\ \textit{Note:} correct because it accurately calculates the complementary probability of no rain as 30\%, demonstrating the fundamental principle in probability theory that the sum of the probabilities of all possible outcomes must equal 1, which is essential for understanding risk assessment in weather forecasting.
\end{enumerate}

\vfill
\noindent\textit{End of Answer Key}
\end{document}